\documentclass[UTF8,a4paper,12pt]{ctexart}

% ========== 数学宏包 ==========
\usepackage{amsmath,amssymb,amsthm}
\usepackage{tikz-cd}
\usepackage{mathtools}

% ========== 页面设置 ==========
\usepackage{geometry}
\geometry{left=2.5cm,right=2.5cm,top=2.5cm,bottom=2.5cm}
\usepackage{fancyhdr}
\pagestyle{fancy}
\fancyhf{}
\fancyhead[L]{\small 范畴深度学习笔记}
\fancyhead[R]{\small 严格 2-范畴}
\fancyfoot[C]{\thepage}
\renewcommand{\headrulewidth}{0.4pt}
\setlength{\headheight}{14pt}

% ========== 定理环境 ==========
\theoremstyle{definition}
\newtheorem{definition}{定义}[section]
\newtheorem{example}{例}[section]
\newtheorem{remark}{注记}[section]
\newtheorem{proposition}{命题}[section]

\theoremstyle{plain}
\newtheorem{axiom}{公理}[section]

% ========== 自定义命令 ==========
\newcommand{\cat}[1]{\mathcal{#1}}
\newcommand{\Ob}{\mathrm{Ob}}
\newcommand{\Hom}{\mathrm{Hom}}
\newcommand{\id}{\mathrm{id}}

\title{\bfseries 严格 2-范畴}
\author{}
\date{\today}

\begin{document}

\maketitle
\thispagestyle{fancy}

\begin{abstract}
本文给出严格 2-范畴（strict 2-category）的精简定义。
严格 2-范畴是范畴概念的「一层提升」：它不仅包含对象和态射，还包含态射之间的态射（称为 2-细胞）。
核心思路是将每个 Hom-集升级为一个 Hom-范畴，并要求水平合成在严格意义下满足范畴公理。
在给出正式定义之前，我们首先回顾必要的前置概念。
\end{abstract}

\tableofcontents
\newpage

% ============================================================
% 第 1 节：前置知识
% ============================================================
\section{前置知识}

\subsection{小范畴与大范畴}

\begin{definition}[小范畴]
若一个范畴 $\cat{C}$ 的对象类 $\Ob(\cat{C})$ 和所有态射构成的类均为集合（而非真类），则称 $\cat{C}$ 为\textbf{小范畴}（small category）。
\end{definition}

\begin{remark}
在严格 2-范畴的定义中，我们要求 $\cat{K}(A,B)$ 为小范畴，这保证了 Hom-范畴自身的良定义性，也使得后续的水平合成双函子在集合论层面无悖论风险。
\end{remark}

\subsection{双函子}

\begin{definition}[双函子]
设 $\cat{A}, \cat{B}, \cat{C}$ 为范畴。一个\textbf{双函子}（bifunctor）
\[
F : \cat{A} \times \cat{B} \to \cat{C}
\]
是满足以下条件的映射：
\begin{enumerate}
  \item 对任意对象 $A \in \cat{A}, B \in \cat{B}$，指定 $\cat{C}$ 中的对象 $F(A,B)$；
  \item 对任意态射 $f : A \to A'$ 于 $\cat{A}$ 及 $g : B \to B'$ 于 $\cat{B}$，指定 $\cat{C}$ 中的态射
  \[
  F(f,g) : F(A,B) \to F(A',B')
  \]
\end{enumerate}
并且满足：
\begin{itemize}
  \item \textbf{对每个变量分别函子性}：固定 $B$，$F(-,B) : \cat{A} \to \cat{C}$ 是函子；固定 $A$，$F(A,-) : \cat{B} \to \cat{C}$ 是函子。
  \item \textbf{兼容性条件}：对任意 $f : A \to A'$, $f' : A' \to A''$ 及 $g : B \to B'$, $g' : B' \to B''$，下图交换：
  \[
  \begin{tikzcd}
    {F(A,B)} \ar[r,"{F(f,g)}"] \ar[d,"{F(\id_A, g' \circ g)}"'] &
    {F(A',B')} \ar[d,"{F(f', g')}"] \\
    {F(A,B'')} \ar[r,"{F(f' \circ f, \id_{B''})}"'] &
    {F(A'',B'')}
  \end{tikzcd}
  \]
  等价地，$F(f',g') \circ F(f,g) = F(f' \circ f, g' \circ g)$。
\end{itemize}
\end{definition}

\begin{remark}
双函子的核心直觉是：它是一个「二元函子」，对每个变量分别函子，且两个变量的作用可以交换（interchange law）。
在 2-范畴中，水平合成正是一个双函子，这一条件强于仅仅在对象层面有复合运算。
\end{remark}

\subsection{范畴的构成要件}

一个范畴 $\cat{C}$ 由以下数据构成：
\begin{itemize}
  \item 对象类 $\Ob(\cat{C})$；
  \item 对任意 $X,Y \in \Ob(\cat{C})$，一个态射集 $\Hom_{\cat{C}}(X,Y)$；
  \item 复合运算 $\circ : \Hom(Y,Z) \times \Hom(X,Y) \to \Hom(X,Z)$；
  \item 恒等态射 $\id_X \in \Hom(X,X)$；
\end{itemize}
满足结合律与单位律。

在严格 2-范畴的语境下，我们将这一构造「提升一层」：对象类变为 0-细胞类，态射集变为 Hom-小范畴，复合运算变为水平合成双函子。

% ============================================================
% 第 2 节：严格 2-范畴的定义
% ============================================================
\section{严格 2-范畴的定义}

\begin{definition}[严格 2-范畴]
一个\textbf{严格 2-范畴}（strict 2-category）$\cat{K}$ 由以下数据构成：
\end{definition}

\subsection{基本数据}

\begin{enumerate}
  \item \textbf{0-细胞类}：一个类 $\Ob(\cat{K})$，其元素称为\textbf{0-细胞}（0-cells），常记作 $A, B, C, \dots$。它们对应于范畴中的「对象」。
  
  \item \textbf{Hom-小范畴}：对任意 $A, B \in \Ob(\cat{K})$，给定一个小范畴 $\cat{K}(A, B)$。
  \begin{itemize}
    \item $\cat{K}(A,B)$ 的对象称为\textbf{1-细胞}（1-cells），即从 $A$ 到 $B$ 的态射；
    \item $\cat{K}(A,B)$ 的态射称为\textbf{2-细胞}（2-cells），即 1-细胞之间的态射。
  \end{itemize}
  
  \item \textbf{水平合成双函子}：对任意 $A, B, C \in \Ob(\cat{K})$，给定一个双函子
  \[
  * : \cat{K}(B,C) \times \cat{K}(A,B) \longrightarrow \cat{K}(A,C)
  \]
  称为\textbf{水平合成}（horizontal composition）。
  
  \item \textbf{恒等 1-细胞}：对任意 $A \in \Ob(\cat{K})$，在 $\cat{K}(A,A)$ 中指定一个对象 $1_A$，称为 $A$ 上的\textbf{恒等 1-细胞}。
\end{enumerate}

\subsection{公理}

上述数据必须满足如下公理：

\begin{axiom}[严格范畴公理]
以 $\Ob(\cat{K})$ 的元素为对象，以 $\cat{K}(A,B)$ 的对象（即 1-细胞）为态射，以双函子 $*$ 在对象上的映射为态射复合，以 $1_A$ 为恒等态射，上述数据严格构成一个范畴。
\end{axiom}

具体而言，该公理等价于以下两个条件的严格成立（以等式而非自然同构的方式）：

\begin{enumerate}
  \item \textbf{结合律（严格）}：对任意 1-细胞 $h : C \to D$, $g : B \to C$, $f : A \to B$，有
  \[
  h * (g * f) = (h * g) * f
  \]
  其中 $*$ 表示双函子在对象（1-细胞）层面的作用。
  
  \item \textbf{单位律（严格）}：对任意 1-细胞 $f : A \to B$，有
  \[
  1_B * f = f = f * 1_A
  \]
\end{enumerate}

\begin{remark}
  这里的「严格」一词正是指结合律和单位律以等式成立，而非仅在同构（2-同构）意义下成立。
  若仅要求 2-同构意义下的结合律与单位律，则得到\textbf{双范畴}（bicategory），这是弱化的 2-范畴概念，在实际应用中更为常见，但其定义也更为复杂。
\end{remark}

% ============================================================
% 第 3 节：图示理解
% ============================================================
\section{图示与直观理解}

\subsection{2-细胞的示意图}

在严格 2-范畴中，数据可以形象地表示为：

\[
\begin{tikzcd}
  A \ar[r,bend left=50,"f"{name=U}] \ar[r,bend right=50,"g"{name=D}] &
  B
  \ar[Rightarrow,from=U,to=D,"\alpha"]
\end{tikzcd}
\]

其中：
\begin{itemize}
  \item $A, B$ 是 0-细胞；
  \item $f, g$ 是 1-细胞（从 $A$ 到 $B$ 的态射）；
  \item $\alpha : f \Rightarrow g$ 是 2-细胞（$f$ 到 $g$ 的态射）。
\end{itemize}

\subsection{水平合成与竖直合成}

严格 2-范畴中有两种合成方式：

\begin{enumerate}
  \item \textbf{竖直合成}（vertical composition）：在同一个 Hom-范畴 $\cat{K}(A,B)$ 内部，对 2-细胞进行合成。这是范畴 $\cat{K}(A,B)$ 自身的态射复合。
  
  \[
  \begin{tikzcd}
    A \ar[r,bend left=60,"f"{name=U}] \ar[r,"h"{name=M,below}] \ar[r,bend right=60,"g"{name=D}] &
    B
    \ar[Rightarrow,from=U,to=M,"\alpha"]
    \ar[Rightarrow,from=M,to=D,"\beta"]
  \end{tikzcd}
  \quad \leadsto \quad
  \begin{tikzcd}
    A \ar[r,bend left=50,"f"{name=U}] \ar[r,bend right=50,"g"{name=D}] &
    B
    \ar[Rightarrow,from=U,to=D,"\beta \circ_v \alpha"]
  \end{tikzcd}
  \]

  \item \textbf{水平合成}（horizontal composition）：通过双函子 $*$，将相邻的 2-细胞并排合成。
  
  \[
  \begin{tikzcd}
    A \ar[r,bend left=50,"f_1"{name=U1}] \ar[r,bend right=50,"g_1"{name=D1}] &
    B \ar[r,bend left=50,"f_2"{name=U2}] \ar[r,bend right=50,"g_2"{name=D2}] &
    C
    \ar[Rightarrow,from=U1,to=D1,"\alpha"]
    \ar[Rightarrow,from=U2,to=D2,"\beta"]
  \end{tikzcd}
  \quad \leadsto \quad
  \begin{tikzcd}
    A \ar[r,bend left=50,"f_2 \circ f_1"{name=U}] \ar[r,bend right=50,"g_2 \circ g_1"{name=D}] &
    C
    \ar[Rightarrow,from=U,to=D,"\beta * \alpha"]
  \end{tikzcd}
  \]
\end{enumerate}

\subsection{交换律}

严格 2-范畴中一个关键的性质是\textbf{交换律}（interchange law），它源于 $*$ 是双函子这一事实：

设 $\alpha : f \Rightarrow f'$, $\alpha' : f' \Rightarrow f''$ 在 $\cat{K}(A,B)$ 中，$\beta : g \Rightarrow g'$, $\beta' : g' \Rightarrow g''$ 在 $\cat{K}(B,C)$ 中。则

\[
(\beta' \circ_v \beta) * (\alpha' \circ_v \alpha) = (\beta' * \alpha') \circ_v (\beta * \alpha)
\]

该等式说明了「先水平后竖直」与「先竖直后水平」的结果一致。这一定理可以从双函子的兼容性条件推导而出，是 2-范畴论中最基本的恒等式之一。

% ============================================================
% 第 4 节：典型例子
% ============================================================
\section{典型例子}

\begin{example}[小范畴的 2-范畴 $\mathbf{Cat}$]
\textbf{$\mathbf{Cat}$}（或更严格地说，全体小范畴构成的 2-范畴）是最基本的例子：
\begin{itemize}
  \item \textbf{0-细胞}：小范畴；
  \item \textbf{1-细胞}：函子 $F : \cat{C} \to \cat{D}$；
  \item \textbf{2-细胞}：自然变换 $\alpha : F \Rightarrow G$；
  \item \textbf{竖直合成}：自然变换的竖直合成（通常的复合）；
  \item \textbf{水平合成}：自然变换的水平合成（Whiskering 与横合成）。
\end{itemize}
在 $\mathbf{Cat}$ 中，结合律与单位律是严格成立的，因此它是一个严格 2-范畴。
\end{example}

\begin{example}[带一个对象的 2-范畴与严格幺半范畴]
若严格 2-范畴 $\cat{K}$ 只有一个 0-细胞 $*$，则：
\begin{itemize}
  \item $\cat{K}(*,*)$ 是一个小范畴，其对象为 1-细胞，态射为 2-细胞；
  \item 水平合成双函子 $* : \cat{K}(*,*) \times \cat{K}(*,*) \to \cat{K}(*,*)$ 赋予 $\cat{K}(*,*)$ 一个严格幺半范畴的结构。
\end{itemize}
反之，任何严格幺半范畴都可视为只有一个对象的严格 2-范畴。这建立了严格幺半范畴与严格 2-范畴之间的等价关系。
\end{example}

% ============================================================
% 第 5 节：与其他概念的关系
% ============================================================
\section{与相关概念的关系}

\subsection{从范畴到 2-范畴}

普通范畴可视为「退化的」2-范畴：若一个严格 2-范畴中，所有 2-细胞均为恒等 2-细胞（即每个 $\cat{K}(A,B)$ 是离散范畴），则其退化为普通范畴。

更精确地说，存在一个忠实的嵌入
\[
\mathbf{Cat} \hookrightarrow \mathbf{2\text{-}Cat}
\]
将每个范畴视为只有平凡 2-细胞的 2-范畴。

\subsection{严格 2-范畴与双范畴}

\begin{itemize}
  \item \textbf{严格 2-范畴}：结合律与单位律以等式严格成立。
  \item \textbf{双范畴（bicategory）}：仅要求结合律与单位律在同构（可逆 2-细胞）意义下成立，且这些同构需满足融贯条件（coherence conditions）。
\end{itemize}

每个严格 2-范畴自然是一个双范畴（取恒等同构即可），但反之不然。在实际范畴论的应用中（如 $\mathbf{Cat}$ 的某些变体、纤维范畴的 2-范畴等），双范畴往往比严格 2-范畴更自然地出现。然而，严格 2-范畴因其简洁的定义，仍是入门学习的理想起点。

\subsection{与 $(n,r)$-范畴层级的关系}

严格 2-范畴是 $(2,2)$-范畴的一个特例——它拥有所有 $k \leq 2$ 维的细胞，且 $k > 2$ 维的细胞都是平凡的。在更广阔的\textbf{高阶范畴论}（higher category theory）图景中，严格 2-范畴是最简单的非平凡高阶结构，为理解弱 $n$-范畴提供了重要的铺垫。

% ============================================================
% 参考文献
% ============================================================
\begin{thebibliography}{9}
\bibitem{leinster2004}
  Tom Leinster,
  \textit{Higher Operads, Higher Categories},
  London Mathematical Society Lecture Note Series 298,
  Cambridge University Press, 2004.

\bibitem{maclane1998}
  Saunders Mac Lane,
  \textit{Categories for the Working Mathematician},
  Second Edition,
  Graduate Texts in Mathematics 5,
  Springer, 1998.

\bibitem{riehl2017}
  Emily Riehl,
  \textit{Category Theory in Context},
  Dover Publications, 2017.

\bibitem{johnson1987}
  Niles Johnson \& Donald Yau,
  \textit{2-Dimensional Categories},
  Oxford University Press, 2021.
\end{thebibliography}

\end{document}